%%%% llars.tex - IJCAI-ECAI 2026 Demo Paper

\typeout{LLARS Demo Paper for IJCAI--ECAI 2026}

\documentclass{article}
\pdfpagewidth=8.5in
\pdfpageheight=11in

\usepackage{ijcai26}

% Packages
\usepackage{times}
\usepackage{soul}
\usepackage{url}
\usepackage[hidelinks]{hyperref}
\usepackage[utf8]{inputenc}
\usepackage[small]{caption}
\usepackage{graphicx}
\usepackage{amsmath}
\usepackage{amsthm}
\usepackage{booktabs}
\usepackage{multirow}
\usepackage{amssymb}
\usepackage[switch]{lineno}

% Comment out this line in the camera-ready submission
\linenumbers

\urlstyle{same}

% PDF Info
\pdfinfo{
/TemplateVersion (IJCAI.2026.0)
}

\title{LLARS: A Collaborative Platform for Interdisciplinary LLM Evaluation}

\author{
Philipp Steigerwald$^{1}$\and
Robin Feldmann$^{1}$\and
Jennifer Burghardt$^{2}$\and
Mara Stieler$^{2}$\And
Eric Rudolph$^{1}$\and
Nico Bienlein$^{1}$\and
Mark Franz$^{1}$\And
Robert Lehmann$^{2}$\and
Jens Albrecht$^{1}$\\
\affiliations
$^{1}$Technische Hochschule N\"urnberg Georg Simon Ohm\\
$^{2}$Institut f\"ur E-Beratung, Technische Hochschule N\"urnberg\\
\emails
philipp.steigerwald@th-nuernberg.de
}

\begin{document}

\maketitle

%==============================================================================
\begin{abstract}
%==============================================================================
Evaluating Large Language Model outputs in specialized domains requires expertise that machine learning researchers often lack, while domain experts lack the technical infrastructure for systematic evaluation.
We present LLARS, an open-source platform enabling interdisciplinary teams to collaboratively evaluate LLM-generated content.
The system provides configurable evaluation workflows including rating scales, bucket-based ranking, and pairwise comparison, with automatic computation of inter-annotator agreement via Krippendorff's~$\alpha$.
LLARS has been deployed in psychosocial online counseling research, supporting studies on automated anonymization and subject line generation.
Beyond evaluation, the platform integrates collaborative prompt engineering, a no-code RAG chatbot builder, and multiple LLM agent reasoning modes.
LLARS is available at \url{https://llars.e-beratungsinstitut.de}.
\end{abstract}

%==============================================================================
\section{Introduction}
%==============================================================================

The deployment of Large Language Models in sensitive domains such as healthcare, legal advice, and psychosocial counseling requires rigorous evaluation by domain experts \cite{chang2024survey}.
However, current evaluation practices face a fundamental gap: the people best qualified to assess output quality often lack the technical infrastructure to conduct systematic evaluations, while ML researchers lack the domain expertise to judge appropriateness.

This challenge is particularly acute in interdisciplinary research settings.
In our work on AI-assisted psychosocial online counseling, counselors must evaluate whether LLM-generated responses are empathetic, clinically appropriate, and safe---judgments that require years of professional training.
Comparing outputs from different prompting strategies, models, and fine-tuning approaches across hundreds of samples becomes cognitively overwhelming without proper tooling.

Existing annotation platforms such as Label Studio \cite{labelstudio2024} and Prodigy \cite{prodigy2024} focus primarily on traditional NLP labeling tasks.
While recent updates have added LLM evaluation templates, these tools were not designed for the specific requirements of expert-driven comparative evaluation with built-in statistical analysis.
Crowdsourcing platforms like Chatbot Arena \cite{zheng2024judging} emphasize preference collection at scale but are unsuitable for sensitive domains requiring specialized expertise.

We present LLARS (LLM Assisted Research System), a platform designed specifically for interdisciplinary LLM evaluation.
LLARS has been successfully deployed in two published studies: evaluating anonymization quality in counseling transcripts \cite{steigerwald2025anonymization} and assessing automated subject line generation for e-mental health \cite{steigerwald2025subjectline}.
The platform is open-source and actively used by research teams combining AI expertise with domain knowledge from psychology and social work.

%==============================================================================
\section{System Overview}
%==============================================================================

LLARS follows a microservices architecture deployed via Docker Compose (see Figure~\ref{fig:architecture}).
The backend provides RESTful APIs and WebSocket connections for real-time collaboration, while the Vue.js frontend offers responsive evaluation interfaces.
Authentication integrates with institutional identity providers via OAuth2/OIDC.
All evaluation data, user annotations, and collaboration artifacts are stored in a relational database with full audit trails.

\begin{figure}[t]
\centering
\fbox{\parbox{0.9\columnwidth}{\centering\vspace{2em}[Architecture Diagram Placeholder]\\\small Flask Backend $\leftrightarrow$ Vue.js Frontend\\MariaDB + ChromaDB + Authentik\vspace{2em}}}
\caption{LLARS system architecture showing the microservices deployment with backend, frontend, and integrated services.}
\label{fig:architecture}
\end{figure}

%==============================================================================
\section{Evaluation Workflows}
%==============================================================================

The core functionality of LLARS centers on structured evaluation scenarios that reduce cognitive load while ensuring data quality.
Table~\ref{tab:eval-modes} summarizes the three evaluation paradigms.

\begin{table}[t]
\centering
\small
\begin{tabular}{@{}lp{5cm}@{}}
\toprule
\textbf{Mode} & \textbf{Description} \\
\midrule
Rating & Evaluators assign scores on configurable scales (e.g., 1--5 stars, Likert). Multiple dimensions per scenario. \\
\addlinespace
Bucket Ranking & Two-phase approach: sort into quality tiers first, then rank within tiers. Reduces absolute judgment burden. \\
\addlinespace
Pairwise & Side-by-side comparison of two outputs with winner selection and optional justification. \\
\bottomrule
\end{tabular}
\caption{Evaluation paradigms supported by LLARS.}
\label{tab:eval-modes}
\end{table}

The bucket-based ranking addresses a key challenge in large-scale evaluation: asking annotators to produce absolute rankings across dozens of samples is cognitively demanding and error-prone.
Instead, evaluators first sort samples into qualitative buckets (e.g., ``excellent'', ``acceptable'', ``poor''), then perform relative ranking only within each bucket.
This two-phase approach significantly reduces the number of required comparisons while maintaining ranking granularity.

LLARS also supports authenticity detection scenarios where evaluators distinguish human-written from AI-generated content, with automatic accuracy computation against ground truth labels.
This functionality supports research on LLM output naturalness and misinformation detection.

All evaluation modes automatically compute Krippendorff's~$\alpha$ \cite{krippendorff2011computing} to quantify inter-annotator agreement.
Results are visualized in real-time dashboards (Figure~\ref{fig:dashboard}), enabling researchers to monitor annotation progress, identify ambiguous samples requiring discussion, and establish evaluation reliability for publication.

\begin{figure}[t]
\centering
\fbox{\parbox{0.9\columnwidth}{\centering\vspace{2em}[Dashboard Screenshot Placeholder]\\\small Real-time agreement metrics\\Progress tracking per evaluator\vspace{2em}}}
\caption{Evaluation dashboard showing Krippendorff's~$\alpha$ and annotation progress across evaluators.}
\label{fig:dashboard}
\end{figure}

%==============================================================================
\section{Collaborative Research Tools}
%==============================================================================

Effective LLM research requires iteration between prompt development and output evaluation.
LLARS integrates these workflows through collaborative editing with real-time synchronization based on Conflict-free Replicated Data Types (CRDTs) \cite{shapiro2011crdt}.

Multiple team members can simultaneously edit prompts, with each user assigned a distinct color for visibility.
A simplified Git-based version control system maintains complete history, showing who edited what and when.
Prompts can be tested against configured LLM backends directly within the interface, with outputs immediately available for comparative evaluation.

The platform follows an \textbf{AI-by-Design} philosophy: at every text input, users can invoke LLM assistance for rephrasing, grammar correction, or content suggestions.
This integration reduces friction between human expertise and AI capabilities throughout the research workflow.

%==============================================================================
\section{Additional Capabilities}
%==============================================================================

Beyond evaluation and prompt engineering, LLARS provides tools for the complete LLM research lifecycle.

The \textbf{RAG Chatbot Wizard} offers a no-code interface for creating domain-specific chatbots.
Users provide a URL, and LLARS automatically crawls the content, generates embeddings, and deploys a Retrieval-Augmented Generation chatbot---enabling domain experts to prototype RAG applications without programming.

LLARS supports multiple \textbf{agent reasoning modes} including Act (direct execution), ReAct \cite{yao2023react} (interleaved reasoning and acting), and Reflect (self-critique before response).
These can be compared side-by-side to evaluate which reasoning strategy performs best for specific domains.

For research teams writing publications, a \textbf{collaborative \LaTeX{} editor} with AI assistance enables real-time co-authoring with the same CRDT-based synchronization and version control used for prompt development.

%==============================================================================
\section{Demonstration}
%==============================================================================

The live demonstration showcases a complete evaluation workflow from the psychosocial counseling domain.
We create an evaluation scenario with rating dimensions for empathy, appropriateness, and helpfulness, then import LLM outputs from different prompting strategies.
Attendees experience the bucket-based ranking interface and observe real-time agreement computation.
Finally, we demonstrate the RAG chatbot wizard by deploying a domain-specific assistant from web content in under two minutes.

%==============================================================================
\section{Related Work}
%==============================================================================

Table~\ref{tab:comparison} compares LLARS with existing platforms.
Label Studio \cite{labelstudio2024} and Argilla \cite{argilla2024} are general-purpose annotation tools that have recently added LLM evaluation features, but lack native support for bucket-based ranking and automatic agreement computation.
Prodigy \cite{prodigy2024} excels at active learning workflows but focuses on individual annotators rather than collaborative team evaluation.
Chatbot Arena \cite{zheng2024judging} enables large-scale preference collection but is designed for crowdsourcing rather than expert evaluation in sensitive domains.

\begin{table}[t]
\centering
\small
\setlength{\tabcolsep}{3pt}
\begin{tabular}{@{}lccccc@{}}
\toprule
& \rotatebox{70}{Bucket Rank.} & \rotatebox{70}{Kripp.~$\alpha$} & \rotatebox{70}{Real-time Collab.} & \rotatebox{70}{Prompt Eng.} & \rotatebox{70}{RAG Builder} \\
\midrule
LLARS & \checkmark & \checkmark & \checkmark & \checkmark & \checkmark \\
Label Studio & -- & -- & \checkmark & (\checkmark) & -- \\
Argilla & -- & -- & \checkmark & -- & -- \\
Prodigy & -- & -- & -- & -- & -- \\
Chatbot Arena & -- & -- & -- & -- & -- \\
\bottomrule
\end{tabular}
\caption{Feature comparison with existing platforms. (\checkmark) indicates partial support.}
\label{tab:comparison}
\end{table}

%==============================================================================
\section{Conclusion}
%==============================================================================

LLARS enables interdisciplinary teams to conduct rigorous LLM evaluations by providing domain experts with professional-grade tooling.
The platform has been validated in psychosocial online counseling research and is designed for extension to other sensitive domains requiring expert judgment.
LLARS is open-source and available at \url{https://llars.e-beratungsinstitut.de}.

\section*{Acknowledgments}
This work was supported by the Bavarian State Ministry of Science and the Arts.

%==============================================================================
% REFERENCES
%==============================================================================
\bibliographystyle{named}
\bibliography{llars}

\end{document}
